% Generated draft: requires mathematical and source review.
\documentclass[11pt]{article}
\usepackage[margin=25mm]{geometry}
\usepackage{fontspec}
\setmainfont{Latin Modern Roman}
\usepackage{amsmath,amssymb,mathtools,unicode-math}
\setmathfont{Latin Modern Math}
\usepackage{enumitem,needspace,xurl,hyperref}
\setlength{\parindent}{0pt}
\setlength{\parskip}{5pt}
\setlength{\emergencystretch}{3em}
\newcommand{\N}{\mathbb{N}}
\newcommand{\Z}{\mathbb{Z}}
\newcommand{\Q}{\mathbb{Q}}
\newcommand{\R}{\mathbb{R}}
\newcommand{\C}{\mathbb{C}}
\newcommand{\F}{\mathbb{F}}
\newcommand{\A}{\mathbb{A}}
\newcommand{\PP}{\mathbb{P}}
\newcommand{\E}{\mathbb{E}}
\newcommand{\eps}{\varepsilon}
\newcommand{\dd}{\,\mathrm{d}}
\begin{document}
\section*{UnsolvedMath ID: 30003232}
\textbf{Source problem code:} catalogue entry
\subsection{Definitions and mathematical statement}
\textbf{Formulation status.} The supplied record does not specify the Diophantine inequalities, their parameters or quantifiers, or the ambient space defining the sets under discussion. Its dated literature assessment describes the extracted item as a continuation labelled conjecture that does not identify a standalone universal assertion. Consequently, the following is a schematic formulation of the recorded questions, not a fully specified conjecture.

Once an ambient metric space \((X,d)\) and the relevant Diophantine conditions have been specified, write \(\operatorname{Bad}\subseteq X\) for the set of vectors satisfying those conditions. For \(A\subseteq X\), let \(\dim_{\mathrm H}A\) denote its Hausdorff dimension with the induced metric. To say that \(A\) has full Hausdorff dimension in \(X\) means
\[
\dim_{\mathrm H}A=\dim_{\mathrm H}X.
\]
Here \(X\), its metric, and the definition of \(\operatorname{Bad}\) are unresolved inputs, not additional hypotheses supplied by the source.

\textbf{Recorded questions.} For a specified notion of badly approximable vectors, determine whether
\[
\dim_{\mathrm H}\operatorname{Bad}=\dim_{\mathrm H}X.
\]
More generally, for a countable family \((\operatorname{Bad}_i)_{i=1}^{\infty}\) of such sets, once their common ambient metric space and their defining conditions are fixed, determine the size and Hausdorff dimension of
\[
B=\bigcap_{i=1}^{\infty}\operatorname{Bad}_i.
\]
The record does not define what additional meaning of ``size'' is intended, nor does it explicitly assert that every such intersection has full Hausdorff dimension. The wording ``for any specified system'' does not provide a class of admissible systems over which a universal assertion can be formulated. No approximation exponents, weights, constants, parameter ranges, boundary conditions, or compatibility assumptions on the countable family are supplied.
\subsection{Short English statement}
The recorded problem asks whether a specified badly approximable set has full Hausdorff dimension and, more generally, how large a countable intersection of such sets is. The defining inequalities and ambient space are missing, so this is only a schematic statement requiring recovery of the surrounding source context. The supplied open label is a dated curation claim, not an independently verified status.
\subsection{Sources}
\begin{itemize}
\item Supplied local UnsolvedMath registry record, permanent ID \(30003232\), title: \emph{Hausdorff Dimension of Intersections of Badly Approximable Sets}; exact source problem code: \texttt{catalogue entry}. Collection: Oberwolfach Reports Open Problems. No individual proposer is identified; the supplied proposed year is 2016.
\item Supplied report citation: \emph{Arbeitsgemeinschaft: Diophantine Approximation, Fractal Geometry and Dynamics (2017)}, Oberwolfach Reports, pp.\ 2749--2792, DOI 10.4171/owr/2016/48. \url{https://doi.org/10.4171/owr/2016/48}. The background attributes the item to workshop open-problem discussions documented in Oberwolfach Reports.
\item Supplied literature source: Badziahin--Velani, ``multiplicative approximation and generalized Cantor sets'' (label as supplied). \url{https://arxiv.org/abs/1007.1848}. The record supplies no more detailed bibliographic information or pinpoint reference for this item.
\item Literature assessment dated exactly \texttt{2026-08-21}: supplied dated curation, not independently verified here. It labels the problem open, describes the extracted item as a continuation labelled conjecture without a standalone universal assertion, and states that the surrounding generalized-Cantor method supplies tools and special cases rather than a complete resolution of the implied problem. The background requests independent verification of time-sensitive status and cited claims.
\item The registry labels the statement \texttt{corrected\_verified} and reports that it was checked against the official Oberwolfach report. This verification claim is supplied metadata, not a source check performed for this draft; it does not resolve the missing formulation details identified in the supplied assessment.
\end{itemize}
\end{document}
